\chapter{Background and Protocol Foundations}
\label{chap:background}

This chapter gives the protocol background needed to understand the implementation work
presented later in the report. It does not introduce new measurements and it does not claim a
new proof of security. Its purpose is to fix the vocabulary, the protocol objects, the role
of the sponsors, and the distinction between the optimistic execution path and the dispute
execution path. This is important because most of the engineering choices made in this
project are only meaningful once the scope of each protocol step is unambiguous.

\section{Fair exchange and optimistic protocols}

A fair exchange protocol allows two parties to exchange valuable objects without trusting
each other. In the setting studied in this project, the vendor \(V\) owns a digital item
\(x\), and the buyer \(B\) is willing to pay a price \(y\) if the item satisfies a public
description predicate. Informally, the desired outcome is:
\[
  B \text{ obtains a valid item } x
  \quad\Longleftrightarrow\quad
  V \text{ obtains the payment } y.
\]
If the exchange cannot complete, neither side should obtain a strictly better position by
deviating from the protocol. A dishonest vendor should not be able to get paid for an invalid
or unavailable item, and a dishonest buyer should not be able to obtain the item while
recovering the payment.

Perfect two-party fair exchange is known to be impossible in the general asynchronous setting
without an additional enforcement mechanism. Blockchain-based protocols use the smart
contract as such a mechanism. The contract is not a trusted party in the traditional sense:
its code is public, deterministic, and executed by the blockchain. It can nevertheless hold
funds, enforce deadlines, verify cryptographic evidence, and redistribute deposits when one
party refuses to continue.

The difficulty is cost. A naive smart-contract solution would put the item, the verification
logic, or a large part of the computation on-chain. This is not viable for large files. The
gas cost of storing data, deploying bytecode, and executing cryptographic computations grows
quickly and can dominate the value of the exchange. The protocol therefore follows an
\emph{optimistic} design: the common case should be cheap, and expensive verification should
only happen if a dispute is raised.

An optimistic fair exchange protocol separates two paths:

\begin{itemize}
  \item the \emph{optimistic path}, where both parties behave as expected and the exchange
  completes after a small number of transactions;
  \item the \emph{dispute path}, where one party claims that the exchange is invalid and the
  contract adjudicates the claim using a compact proof.
\end{itemize}

This separation is central to \sox{}. The protocol accepts that a dispute may be expensive,
but it tries to make disputes rare, logarithmic in the size of the computation trace, and
economically financed by sponsors and deposits.

\section{Overview of the SOX protocol}

The Sponsored Fair Exchange protocol, abbreviated \sox{}, is described in
\cite{sox2026}. The implementation inherited at the beginning of this project follows the
same high-level architecture and was documented in the previous implementation report
\cite{hana2025}.

The digital item is not sent directly to the buyer in clear. Instead, the vendor encrypts
the item and publishes an encrypted object \(ct\). The buyer can download \(ct\), but cannot
recover \(x\) until the vendor reveals the decryption key. Before the payment is released,
the vendor also commits to the data and to the computation that should prove that the
plaintext satisfies the expected description.

The protocol is built around the following public and authenticated objects:

\begin{itemize}
  \item the ciphertext \(ct\), which is the encrypted representation of the item;
  \item the description \(\mathsf{desc}\), for example \(\mathsf{SHA256}\) in the main
  specialized case studied in this project;
  \item the expected digest or description output \(d\);
  \item a circuit \(C\) that evaluates whether decrypting \(ct\) gives an item satisfying
  \(\mathsf{desc}(x)=d\);
  \item cryptographic commitments to \(ct\), to the circuit, and to intermediate values used
  during dispute resolution.
\end{itemize}

The smart contract does not store the whole file and does not execute the full circuit in the
optimistic case. It only stores short commitments and enough economic state to decide who
should receive the funds. If the buyer accepts the decrypted item, the contract releases the
payment to the vendor. If the buyer complains, the dispute mechanism narrows the disagreement
to a single circuit position and verifies a local computation step on-chain.

The important architectural idea is that the contract only needs to verify a small local
claim after a logarithmic search. The expensive work of encrypting, constructing the circuit,
computing intermediate values, and preparing Merkle proofs remains off-chain. The on-chain
part checks enough authenticated data to make lying economically and cryptographically
detectable.

\section{Roles, sponsors, deposits, and incentives}

The protocol distinguishes between the parties who exchange the item and the parties who pay
gas or finance parts of the protocol. In the most general form, there are five roles:

\begin{table}[htbp]
  \centering
  \begin{tabularx}{\linewidth}{>{\raggedright\arraybackslash}p{0.16\linewidth}>{\raggedright\arraybackslash}X}
    \toprule
    Role & Function in the protocol \\
    \midrule
    \(V\) &
    Vendor. Owns the item \(x\), prepares the encrypted item and the precontract data,
    reveals the decryption key, and defends the correctness of the exchange during a
    dispute. \\
    \(B\) &
    Buyer. Accepts the precontract, pays the price, decrypts the ciphertext after receiving
    the key, and raises a complaint if the decrypted item does not satisfy the description. \\
    \(S\) &
    Optimistic sponsor. In the sponsored version, this role finances or deposits funds for
    the optimistic part of the protocol. \\
    \(S_B\) &
    Buyer-side dispute sponsor. This role finances the buyer side of the dispute when the
    buyer claims that the exchange is invalid. \\
    \(S_V\) &
    Vendor-side dispute sponsor. This role finances the vendor side of the dispute when the
    vendor defends the validity of the exchange. \\
    \bottomrule
  \end{tabularx}
  \caption{Protocol roles used by \sox{}.}
  \label{tab:background-roles}
\end{table}

The motivation for introducing sponsors is practical. A buyer or vendor may not want to hold
gas funds, or may not want to expose themselves to the full execution cost of a dispute.
Sponsors can cover these costs in exchange for deposits, tips, or reimbursement rules
encoded in the contract. The exact economic values are implementation parameters, but the
logic is protocol-driven: if a party causes an unnecessary or losing dispute, the deposits
should compensate the party who paid for the successful execution.

Several roles may coincide. For example, \(S=B\) means that the buyer is also the optimistic
sponsor, \(S=V\) means that the vendor is also the optimistic sponsor, \(S_B=B\) means that
the buyer sponsors its own side of the dispute, and \(S_V=V\) means that the vendor sponsors
its own side. These equalities are not cosmetic. They can remove transactions, deposits, or
repeated dispute iterations. A large part of this project studies which coincidences have a
real gas impact and which ones are mostly semantic.

One variant is especially important: the no \(S\)-deposit mode, also called \simsox{} in the
measurements. In that mode there is no initial sponsor deposit by \(S\) in the optimistic
part. The parties pay for their own optimistic transactions, so the contract no longer needs
to carry reimbursement logic for an external optimistic sponsor. This simplifies the
optimistic contract and is one of the main reasons for the lower marginal cost measured in
the final implementation.

\section{Optimistic path of the protocol}

The optimistic path is the path followed when the vendor provides a valid item and the buyer
accepts it. In the inherited implementation and in the specialized variants developed during
this project, the path can be summarized as follows.

\begin{enumerate}
  \item \textbf{Precontract generation.} The vendor prepares the encrypted file, the public
  description data, and the commitments required by the protocol. This computation is
  off-chain and may be expensive for large files.

  \item \textbf{Contract creation or initialization.} The optimistic contract for the
  exchange is deployed or created through the selected deployment mechanism. In the initial
  implementation this was essentially an account-per-exchange design. In the final
  specialized implementation, several modes use a factory and smaller dedicated contracts.

  \item \textbf{Buyer payment.} The buyer sends the payment to the optimistic contract. The
  contract escrows the payment instead of transferring it immediately to the vendor.

  \item \textbf{Key release.} The vendor reveals the symmetric key. The buyer can now decrypt
  the ciphertext and check the item locally.

  \item \textbf{Completion or complaint.} If the buyer is satisfied, the buyer completes the
  exchange and the contract releases the payment to the vendor. If the buyer is not
  satisfied, the protocol leaves the optimistic path and enters the dispute path.
\end{enumerate}

In the Solidity implementation, the optimistic contract states reflect this structure. The
contract waits for the buyer payment, then waits for the vendor key, then waits for either
buyer completion or dispute preparation. The exact state names vary slightly across the
initial, monolithic, and specialized contracts, but the core sequence remains:
\[
  \texttt{WaitPayment}
  \rightarrow
  \texttt{WaitKey}
  \rightarrow
  \texttt{WaitSB}
  \rightarrow
  \texttt{WaitSV} \text{ or } \texttt{End}.
\]
The state \texttt{End} corresponds to a completed or otherwise finalized exchange.

The gas objective for this path is strict: the optimistic path is the common case, so it must
be as cheap as possible. Any branch that only matters in rare dispute scenarios should not be
paid by every honest exchange if it can be moved to a specialized contract or delayed until a
dispute actually occurs.

\section{Dispute path of the protocol}

A dispute starts when the buyer claims that the decrypted item is invalid. The purpose of the
dispute is not to recompute the whole circuit on-chain. Instead, the parties conduct an
interactive search over the computation trace until the disagreement is localized. Then the
contract verifies one local computation step with Merkle proofs and the relevant gate
evaluation.

At a high level, the dispute mechanism has three phases.

\paragraph{Dispute initialization.}
The buyer announces dissatisfaction and the dispute infrastructure is initialized. In the
implementation, the operation usually called \texttt{triggerDispute} corresponds to the
paper's Step 6. It is important to note that in the measured implementation this action can
include deployment or creation of the dispute contract. Therefore, its gas cost is not merely
the cost of changing an optimistic state variable.

\paragraph{Binary search over the computation trace.}
The buyer and the vendor narrow the interval in which they disagree. The implementation uses
\texttt{respondChallenge} for the buyer response and \texttt{giveOpinion} for the vendor's
opinion. These correspond respectively to Step 7b and Step 7c in the protocol mapping used
throughout this report. The number of rounds is logarithmic in the number of gates. For
example, if the circuit has \(N\) gates, the search requires roughly \(\lceil \log_2 N\rceil\)
rounds before it reaches a single local position.

\paragraph{Local verification and finalization.}
Once the relevant position is identified, the vendor submits the local data required to
verify one gate or one boundary condition. In the implementation this is the
\texttt{submitCommitment} family of calls, corresponding to Step 8. The contract checks
membership proofs against the commitments, evaluates the local gate, and then finalizes or
continues according to the protocol logic. The final transition is referred to as
\texttt{finalize} in the benchmarks and corresponds to Step 9.

\begin{table}[htbp]
  \centering
  \begin{tabularx}{\linewidth}{>{\raggedright\arraybackslash}p{0.34\linewidth}>{\raggedright\arraybackslash}p{0.13\linewidth}>{\raggedright\arraybackslash}X}
    \toprule
    Implementation name & Paper step & Purpose \\
    \midrule
    \texttt{triggerDispute} &
    Step 6 &
    Leaves the optimistic path and initializes the dispute game. In the measured contracts,
    this may include dispute deployment. \\
    \texttt{respondChallenge} &
    Step 7b &
    Buyer response during the binary search over the computation trace. \\
    \texttt{giveOpinion} &
    Step 7c &
    Vendor opinion on the buyer response; updates the search interval. \\
    \texttt{submitCommitment}\newline
    \texttt{submitCommitmentLeft}\newline
    \texttt{submitCommitmentRight} &
    Step 8 &
    Local verification of a circuit gate, boundary value, or final comparison using
    authenticated data. \\
    \texttt{finalize} &
    Step 9 &
    Ends the current dispute execution or triggers the next protocol consequence according
    to the sponsor configuration. \\
    \bottomrule
  \end{tabularx}
  \caption{Mapping between implementation terminology and protocol steps.}
  \label{tab:background-step-mapping}
\end{table}

This mapping is used later in the experimental chapter. Without it, gas tables are easy to
misread. For example, a line labelled ``dispute cost'' may include only Step 8, the whole
post-trigger dispute, or also the deployment cost of the dispute contract. This project
therefore reports these scopes separately.

\section[Circuit commitments: hct, hcircuit, and hpre]{Circuit commitments: \(h_{ct}\), \(h_{circuit}\), and \(h_{pre}\)}

The dispute mechanism relies on commitments that allow the contract to verify small pieces of
a large off-chain computation. The implementation uses Merkle-style commitments over
fixed-size objects. The precise encoding is discussed later in the report because it also
raises security and ambiguity questions, but the conceptual role of each commitment is as
follows.

\begin{table}[htbp]
  \centering
  \begin{tabularx}{\linewidth}{>{\raggedright\arraybackslash}p{0.18\linewidth}>{\raggedright\arraybackslash}X}
    \toprule
    Commitment & Role in the protocol \\
    \midrule
    \(h_{ct}\) &
    Commitment to the ciphertext blocks. During Step 8, a party can reveal only the block
    needed for a local gate evaluation, together with a membership proof. \\
    \(h_{circuit}\) &
    Commitment to the circuit description. In the generic mode, the contract does not know
    the whole circuit, so the vendor must authenticate the gate being evaluated. \\
    \(h_{pre}\) &
    Commitment or authenticated value associated with intermediate computation values used
    during the dispute search. It lets the parties argue about a particular prefix or local
    position without publishing the full trace on-chain. \\
    \bottomrule
  \end{tabularx}
  \caption{Commitments used to keep the dispute verification local.}
  \label{tab:background-commitments}
\end{table}

The circuit used by the inherited implementation is encoded as a sequence of fixed-size
64-byte gates. A gate may represent, for example, an AES-CTR decryption block, a SHA-256
compression step, a constant, an XOR operation, or a comparison. The contract can evaluate a
single gate if it receives the gate encoding, the necessary input values, and proofs that
those inputs are consistent with the committed ciphertext or intermediate trace.

The role of \(h_{circuit}\) is especially important for understanding the hardcoded
\(\mathsf{SHA256}\) optimization. In the generic mode, the contract must be convinced that
the gate submitted in Step 8 is really the gate committed to by the precontract. This is why
the generic Step 8 includes a proof for the circuit commitment, denoted \(\pi_1\) in the
project discussions. In the hardcoded \(\mathsf{SHA256}\) mode, the contract can reconstruct
the expected gate from public parameters such as the file length, the description hash, and
the gate index. In that specialized case, the proof \(\pi_1\) is no longer needed for the
same purpose because the circuit structure is no longer supplied by the vendor.

This optimization does not remove all Step 8 costs. The contract still has to verify the
relevant ciphertext or intermediate-value proofs and still has to evaluate the local gate.
It only removes the generic circuit-authentication part of Step 8. This distinction explains
why the local hardcoded gain can be visible on the Step 8 subcomponent while being diluted in
the total cost of a complete dispute.

\section{ERC-4337 and transaction sponsorship}

The implementation also contains a transaction-sponsorship layer inspired by ERC-4337
account abstraction. Instead of requiring each user to send every transaction directly from a
funded externally owned account, the application can package calls as user operations and
send them through a bundler and an EntryPoint contract.

This layer is important for usability because it makes it possible to demonstrate sponsored
execution: a participant can trigger protocol actions without manually managing gas in the
same way as a regular Ethereum transaction. It is also important experimentally because it
introduces infrastructure that is not part of the mathematical core of the protocol. A
bundler, an EntryPoint, and local deployment scripts can affect latency, failure modes, and
the way gas is paid.

For this reason, the report distinguishes two notions of sponsorship:

\begin{itemize}
  \item \emph{protocol sponsorship}, which is the role of \(S\), \(S_B\), and \(S_V\) in the
  \sox{} protocol and determines deposits, reimbursements, and dispute incentives;
  \item \emph{transaction sponsorship}, which is the implementation mechanism that submits
  calls through an account-abstraction stack.
\end{itemize}

These notions are related but not identical. The no \(S\)-deposit mode simplifies protocol
sponsorship in the optimistic phase. It does not automatically remove every piece of
transaction infrastructure from the application. Conversely, an ERC-4337 user operation can
be used to submit a call whose economic meaning is still governed by the protocol roles and
deposits. Keeping this distinction explicit was necessary to avoid confusing gas paid to
execute Ethereum transactions with the protocol deposits used to compensate sponsors.

\section{Security assumptions and fairness goals}

The security goal of \sox{} is best understood operationally. If the vendor behaves honestly,
publishes a valid encrypted item, and reveals the correct key, then the vendor should receive
the payment unless the buyer can provide valid evidence of a problem. If the buyer behaves
honestly and the decrypted item is invalid, unavailable, or inconsistent with the
precontract, then the buyer should be able to prevent the vendor from keeping the payment by
following the dispute procedure. Sponsors should be reimbursed or penalized according to the
rules encoded in the contract.

The implementation relies on the following assumptions.

\begin{itemize}
  \item The blockchain executes the deployed smart contracts deterministically and makes their
  state publicly available.
  \item Deadlines are long enough for honest parties or their sponsors to submit the required
  transactions.
  \item The hash functions used for commitments and Merkle proofs are collision resistant for
  the encoded objects being committed.
  \item The encryption scheme keeps \(x\) confidential before the key is revealed.
  \item The off-chain computation used to produce ciphertexts, circuits, intermediate values,
  and proofs is consistent with the on-chain verifier.
  \item The data needed for a dispute, especially ciphertext blocks and proof artifacts,
  remains available to the party that must use it.
  \item The transaction-sponsorship infrastructure, when used, is live enough not to censor
  honest protocol actions indefinitely.
\end{itemize}

These assumptions also show the limits of the current prototype. The smart contract can
check authenticated local claims, but it cannot force an unavailable backend to serve data,
cannot make a bundler include a transaction, and cannot compensate for ambiguous low-level
encodings if such encodings allow two interpretations of the same byte string. The latter
point became important during the project when reviewing the inherited gate encoding and
Merkle construction. The specialized gas optimizations developed in this project do not
depend on creating new encodings, but a production-grade version should make those encodings
canonical and domain-separated.

The rest of the report builds on these foundations. Chapter~3 describes the inherited
implementation and its execution path. Later chapters then explain how the protocol variants
were isolated, how the final specialized contracts were implemented, and how the gas and
off-chain timings were measured.
